Este capítulo descreve o HyBolt, implementação em Python do HySAIL com as propostas do Capítulo~4. São apresentados o mapeamento entre as fases do protocolo e os comandos da ferramenta, o fluxo de cada fase, a organização do código, os metadados mantidos pelo usuário, os problemas observados ao realizar a substituição de forma direta e as simplificações adotadas em relação a um ambiente de nuvem real.

\section{Visão Geral}

O HyBolt é distribuído como uma ferramenta de linha de comando cujos comandos correspondem às fases do HySAIL, conforme a Tabela~\ref{tab:fases_comandos}. A codificação RaptorQ utiliza a biblioteca \texttt{raptorq}, que implementa a RFC~6330 \cite{luby2011rfc6330}, e a aritmética sobre $GF(2)[x]$, incluindo as divisões polinomiais, o produto sem transporte (\emph{carry-less}) e o teste de primitividade, foi implementada no próprio projeto.

\begin{table}[h]
	\centering
	\caption{Correspondência entre as fases do HySAIL e os comandos do HyBolt}
	\label{tab:fases_comandos}
	\begin{tabular}{ll}
		\hline
		\textbf{Fase do HySAIL} & \textbf{Comando do HyBolt} \\
		\hline
		Codificação e Distribuição & \texttt{hybolt encode} \\
		Auditoria & \texttt{hybolt audit} \\
		Redistribuição & \texttt{hybolt audit --redistribute} \\
		Decodificação & \texttt{hybolt decode} \\
		\hline
	\end{tabular}
\end{table}

Dois esquemas de prova de recuperabilidade estão disponíveis e são escolhidos na codificação pela opção \texttt{--proof-scheme}: o esquema limitado, que corresponde às adaptações mínimas das Seções~\ref{sec:impacto_mac} a~\ref{sec:redistribuicao}, e o esquema compacto da Seção~\ref{sec:esquema_compacto}, adotado como padrão. A ocultação da estrutura do código (Seção~\ref{sec:ocultacao}) é aplicada por padrão a ambos e pode ser desabilitada pela opção \texttt{--no-seal-packets}.

\section{Fluxo das Fases}

\subsection{Codificação e Distribuição}
O arquivo é lido em \emph{chunks} independentes de, por padrão, $\lambda \cdot 512$ bytes. Cada \emph{chunk} é dividido em $K$ símbolos de origem, cujo tamanho padrão corresponde a 10\% do \emph{chunk}, e codificado em $K + \lceil \varepsilon K \rceil$ pacotes RaptorQ, com $\varepsilon$ configurável pela opção \texttt{--redundancy}. Em seguida, cada pacote é cifrado, e sobre o pacote cifrado são calculados o \emph{hash} SHA-256 e o material de prova: no esquema limitado, os $l = 40$ \acrshort{MAC}s da Equação~\eqref{eq:mac_por_pacote}, mantidos pelo usuário, com um conjunto de polinômios primitivos de grau $\lambda = 127$ compartilhado por todo o arquivo; no esquema compacto, a etiqueta da Equação~\eqref{eq:etiqueta_compacta}, enviada ao provedor junto ao pacote.

Os pacotes de cada \emph{chunk} são embaralhados e atribuídos aos provedores de forma circular segundo a posição no embaralhamento, o que torna a alocação aleatória, como no HySAIL, e mantém a carga equilibrada. Cada provedor armazena apenas o pacote cifrado e, no esquema compacto, sua etiqueta, sob um nome pseudoaleatório. Ao final, o usuário grava os metadados descritos na Seção~\ref{sec:metadados} e descarta os pacotes.

\subsection{Auditoria}
Cada execução de \texttt{hybolt audit} realiza uma rodada de auditoria sobre uma amostra aleatória dos pacotes de cada provedor, cujo tamanho é definido pela opção \texttt{--sample-size} (por padrão, todos os pacotes do provedor). No esquema compacto, são sorteados os coeficientes $\nu_\eta$ da Equação~\eqref{eq:resposta_compacta}, e a resposta é verificada pela Equação~\eqref{eq:verificacao_compacta} sem que o polinômio secreto seja transmitido. No esquema limitado, um polinômio ainda não revelado é sorteado e registrado nos metadados antes de qualquer comunicação com os provedores, e são sorteados os expoentes secretos $k_\eta$ da Equação~\eqref{eq:desafio_ponderado}. Em ambos os casos, se a resposta de um provedor divergir da esperada, os pacotes da amostra são desafiados individualmente para localizar os corrompidos, e a ausência de um pacote ou de sua etiqueta é tratada como corrupção. O relatório informa $\tau_{provider}$ de cada provedor, comparado ao limite \texttt{--tau-maximum}, e, para cada \emph{chunk} afetado, se o número de pacotes corrompidos ainda é tolerado por seus pacotes de reparo. O comando termina com código de saída diferente de zero quando encontra corrupção, o que permite sua execução periódica por agendadores.

\subsection{Redistribuição}
Com a opção \texttt{--redistribute}, os pacotes corrompidos ou perdidos encontrados na auditoria são substituídos conforme a Seção~\ref{sec:redistribuicao}: o \emph{chunk} afetado é decodificado a partir dos pacotes cujo SHA-256 confere e recodificado com os mesmos parâmetros, e, para cada pacote perdido, é gerado um novo símbolo de reparo com \acrshort{ESI} inédito, que é cifrado, recebe etiqueta ou \acrshort{MAC}s e é gravado no provedor menos carregado do \emph{chunk}. A opção \texttt{--extra-repair} acrescenta símbolos de reparo adicionais a cada \emph{chunk} reparado. As cópias corrompidas são removidas e os metadados, incluindo o contador de \acrshort{ESI}s, a margem de cada \emph{chunk} e o \emph{hash} raiz que resume a localização e o \emph{hash} de todos os pacotes, são atualizados. Quando um provedor deixa de constar da lista de provedores, a auditoria considera perdidos todos os seus pacotes, e a decodificação prossegue com os pacotes dos demais provedores. Uma nova auditoria após a Redistribuição confirma a restauração.

\subsection{Decodificação}
Na decodificação, os pacotes de cada \emph{chunk} são transferidos, verificados, decifrados e entregues ao decodificador RaptorQ até que os dados sejam reconstruídos. No esquema compacto, cada pacote transferido é verificado contra a etiqueta armazenada pelo provedor, sem consumir desafios. No esquema limitado, o provedor é desafiado antes de cada transferência com um único polinômio sorteado para toda a sessão. Pacotes que não passam na verificação ou cujo \emph{hash} SHA-256 não confere são descartados, e a decodificação prossegue com os pacotes redundantes. A verificação pode ser desabilitada pela opção \texttt{--no-retrievability-proof}, mantendo-se a verificação por SHA-256.

\section{Organização do Código}

A prova de recuperabilidade, a auditoria e a decodificação foram organizadas segundo o padrão de projeto \emph{Strategy}, de modo que variantes possam ser substituídas sem alterar o fluxo principal. As principais classes estão representadas na Figura~\ref{fig:estrutura_codigo} e descritas a seguir:

\begin{figure}[htbp]
	\centering
	\begin{adjustbox}{max width=\textwidth}
		\begin{tikzpicture}[
		interface/.style={draw=serieA!80!black, fill=serieA!25, font=\footnotesize, align=center, inner sep=4pt, text width=42mm},
		concreta/.style={draw=#1!80!black, fill=#1!10, font=\scriptsize\ttfamily, align=center, inner sep=3pt, text width=42mm},
		herda/.style={-{Triangle[open, length=2.5mm, width=2.5mm]}, draw=tintaSec, thick},
		barramento/.style={draw=tintaSec, thick},
		usa/.style={->, dashed, draw=tintaSec},
		node distance=3mm and 6mm]

	% Prova de recuperabilidade
	\node[interface] (rp) {$\langle\langle$interface$\rangle\rangle$\\\texttt{RetrievabilityProof}};
	\node[concreta=serieB, below=7mm of rp] (poly) {PolynomialRetrievabilityProof};
	\node[concreta=serieB, below=of poly] (compact) {CompactRetrievabilityProof};
	\node[concreta=serieB, below=of compact] (disabled) {DisabledRetrievabilityProof};
	\draw[herda] ($(disabled.west)+(-4mm,0)$) -- ($(rp.west)+(-4mm,0)$) -- (rp.west);
	\foreach \c in {poly,compact,disabled} \draw[barramento] (\c.west) -- ++(-4mm,0);
	\node[concreta=serieD, below=7mm of disabled] (ledger) {ChallengeLedger};
	\node[concreta=serieD, below=of ledger] (keys) {CompactKeys};
	\draw[usa] (ledger.east) -- ++(5mm,0) |- (poly.east);
	\draw[usa] (keys.east) -- ++(3mm,0) |- (compact.east);

	% Auditoria
	\node[interface, right=12mm of rp] (ac) {$\langle\langle$interface$\rangle\rangle$\\\texttt{AuditChallenge}};
	\node[concreta=serieB, below=7mm of ac] (bounded) {BoundedAuditChallenge};
	\node[concreta=serieB, below=of bounded] (compactac) {CompactAuditChallenge};
	\draw[herda] ($(compactac.west)+(-4mm,0)$) -- ($(ac.west)+(-4mm,0)$) -- (ac.west);
	\foreach \c in {bounded,compactac} \draw[barramento] (\c.west) -- ++(-4mm,0);
	\node[concreta=serieC, below=7mm of compactac] (audit) {Audit};
	\node[concreta=serieC, below=of audit] (redist) {Redistribution};
	\draw[usa] (audit.east) -- ++(4mm,0) |- (compactac.east);
	\draw[usa] ($(audit.east)+(4mm,0)$) |- (bounded.east);

	% Decodificação
	\node[interface, right=12mm of ac] (ds) {$\langle\langle$interface$\rangle\rangle$\\\texttt{DecodingStrategy}};
	\node[concreta=serieB, below=7mm of ds] (rq) {Raptor10DecodingStrategy};
	\node[concreta=serieB, below=of rq] (lt) {LtDecodingStrategy};
	\draw[herda] ($(lt.west)+(-4mm,0)$) -- ($(ds.west)+(-4mm,0)$) -- (ds.west);
	\foreach \c in {rq,lt} \draw[barramento] (\c.west) -- ++(-4mm,0);
	\node[concreta=serieE, below=7mm of lt] (sealing) {SealingKeys};
	\node[concreta=serieE, below=of sealing] (server) {Server};

	\node[rotulo, above=2mm of rp] {prova de recuperabilidade};
	\node[rotulo, above=2mm of ac] {auditoria e redistribuição};
	\node[rotulo, above=2mm of ds] {decodificação e armazenamento};
\end{tikzpicture}

	\end{adjustbox}
	\caption{Principais classes do HyBolt. As setas com triângulo indicam a realização de uma interface, e as setas tracejadas, dependências.}
	\label{fig:estrutura_codigo}
\end{figure}

\begin{itemize}
	\item \textit{Prova de recuperabilidade:} a interface \texttt{RetrievabilityProof} define o material de prova gerado na codificação (\acrshort{MAC}s do usuário ou etiquetas dos provedores) e as verificações antes e depois de cada transferência. É especializada por \texttt{PolynomialRetrievabilityProof}, que realiza o esquema limitado, por \texttt{CompactRetrievabilityProof}, que realiza o esquema compacto a partir das chaves \texttt{CompactKeys}, e por \texttt{DisabledRetrievabilityProof}, um objeto nulo que permite desabilitar a prova, por exemplo, para medir o custo do restante do protocolo.
	\item \textit{Auditoria:} a interface \texttt{AuditChallenge} representa o desafio de uma rodada e é especializada por \texttt{BoundedAuditChallenge} e \texttt{CompactAuditChallenge}. A classe \texttt{Audit} concentra a amostragem, a localização dos pacotes corrompidos e a avaliação dos limites $\tau$, comuns aos dois esquemas, e a classe \texttt{ChallengeLedger} controla o uso único dos polinômios do esquema limitado.
	\item \textit{Decodificação:} a interface \texttt{DecodingStrategy} concentra o procedimento comum de verificação, transferência e conferência do \emph{hash} de cada pacote, e é especializada por \texttt{RaptorQDecodingStrategy}, utilizada na proposta, e por \texttt{LtDecodingStrategy}, que implementa a decodificação por \emph{peeling} e o cálculo do \acrshort{MAC} esperado pela combinação homomórfica dos \acrshort{MAC}s dos blocos de mensagem, como no HySAIL original. A estratégia é escolhida pelo código registrado nos metadados de cada \emph{chunk}.
	\item \textit{Ocultação da estrutura:} a classe \texttt{SealingKeys} cifra os pacotes e deriva seus nomes, e a classe \texttt{Redistribution} implementa a fase correspondente para ambos os esquemas.
\end{itemize}

\section{Metadados do Usuário}
\label{sec:metadados}

O usuário mantém um arquivo de metadados por arquivo armazenado, que constitui todo o estado necessário às demais fases e deve permanecer sob seu controle, pois contém as chaves secretas. Seu conteúdo é resumido na Tabela~\ref{tab:metadados}.

\begin{table}[h]
	\centering
	\caption{Conteúdo dos metadados mantidos pelo usuário}
	\label{tab:metadados}
	\begin{tabular}{lll}
		\hline
		\textbf{Escopo} & \textbf{Informação} & \textbf{Esquema} \\
		\hline
		Arquivo & nome e SHA-256 do arquivo original & ambos \\
		Arquivo & chaves de cifragem e de nomeação dos pacotes & ambos \\
		Arquivo & \emph{hash} raiz da localização dos pacotes & ambos \\
		Arquivo & polinômio secreto, chave da PRF e $id_F$ & compacto \\
		Arquivo & polinômios de desafio e índices já revelados & limitado \\
		\emph{Chunk} & tamanho, parâmetros RaptorQ e pacotes de reparo & ambos \\
		Pacote & \acrshort{ESI}, provedor, arquivo e SHA-256 & ambos \\
		Pacote & um \acrshort{MAC} por polinômio de desafio & limitado \\
		\hline
	\end{tabular}
\end{table}

\section{Problemas Observados na Substituição Direta}
\label{sec:problemas_observados}

Para quantificar os problemas descritos na Seção~\ref{sec:problemas_substituicao}, o arquivo de 1~MiB do conjunto de exemplos foi codificado com os dois esquemas, com os parâmetros padrão e três provedores, e em seguida auditado, com 50 pacotes por provedor, e decodificado. O arquivo foi dividido em 17 \emph{chunks}, cada um com $K = 11$ pacotes de origem e $r = 3$ pacotes de reparo de 6.500 bytes, totalizando 238 pacotes. Os resultados, obtidos em uma única execução por esquema com Python~3.13 e a versão 2.0.0 da biblioteca \texttt{raptorq}, são apresentados na Tabela~\ref{tab:comparacao_esquemas} e têm caráter indicativo.

\begin{table}[h]
	\centering
	\caption{Comparação entre o esquema limitado e o esquema compacto para um arquivo de 1~MiB}
	\label{tab:comparacao_esquemas}
	\begin{tabular}{lrr}
		\hline
		\textbf{Medida} & \textbf{Limitado} & \textbf{Compacto} \\
		\hline
		Tempo de codificação (s) & 97,4 & 0,72 \\
		Metadados do usuário (bytes) & 1.608.436 & 36.550 \\
		Metadados em relação ao arquivo & 153,4\% & 3,5\% \\
		Auditorias disponíveis & 40 & ilimitadas \\
		Rodada de auditoria (s) & 1,48 & 0,05 \\
		Resposta por provedor (bytes) & 16 & $\approx$ 6.550 \\
		Tempo de decodificação (s) & 1,87 & 0,40 \\
		\hline
	\end{tabular}
\end{table}

A substituição direta, que corresponde ao esquema limitado, confirma os itens 1 e 2 da Seção~\ref{sec:problemas_substituicao}. Os metadados superaram o próprio arquivo, pois cada um dos $238 \cdot 40$ \acrshort{MAC}s é armazenado com um byte por coeficiente; mesmo com o empacotamento em bits, ocupariam cerca de 152~KB, ou 14,5\% do arquivo, valor coerente com a estimativa $l\lambda/B$ acrescida dos pacotes de reparo. A codificação foi cerca de 135 vezes mais lenta que no esquema compacto, por exigir $l = 40$ divisões polinomiais por pacote em vez de uma, e o arquivo admite apenas 40 auditorias ou decodificações com prova antes de ser recodificado. No esquema compacto, os metadados restantes correspondem à localização e ao \emph{hash} de cada pacote; as chaves da prova ocupam 64 bytes. Em contrapartida, a resposta de cada provedor passa a ter o tamanho aproximado de um pacote, em vez de 16 bytes.

O item 3, o ataque dirigido a um \emph{chunk}, também se confirma nesse cenário. Sem a ocultação da estrutura, corromper os 4 pacotes de um mesmo \emph{chunk}, apenas 1,7\% do total, torna o arquivo irrecuperável. Para um arquivo de 1~GiB codificado com os mesmos parâmetros, com aproximadamente 243.700 pacotes, os mesmos 4 pacotes representariam 0,0016\% do total, e uma auditoria com 50 pacotes por provedor em três provedores os detectaria com probabilidade de apenas 0,25\%. Com a ocultação, os nomes dos arquivos, o cabeçalho RaptorQ e os símbolos de origem deixam de ser visíveis aos provedores, como verificado nos testes automatizados, e o adversário perde a capacidade de identificar os pacotes de um mesmo \emph{chunk}.

\section{Simplificações}

A implementação tem caráter de prova de conceito e adota as seguintes simplificações em relação ao modelo do HySAIL:

\begin{itemize}
	\item \textit{Provedores locais:} cada provedor é representado por um diretório no sistema de arquivos local, e as respostas aos desafios são calculadas no mesmo processo do usuário. O canal privado e autenticado entre usuário e provedores pressuposto pelo HySAIL \cite{amaral2019hysail} não é implementado, embora a interface dos provedores, restrita a armazenar, remover, responder desafios e transferir pacotes e etiquetas, delimite o ponto em que uma comunicação remota seria introduzida.
	\item \textit{Ocultação parcial:} os pacotes do último \emph{chunk} podem ter tamanho diferente dos demais, e os pacotes são enviados aos provedores \emph{chunk} a \emph{chunk}, de modo que o tamanho e a ordem de chegada ainda revelam parte da estrutura do código (Seção~\ref{sec:ocultacao}).
	\item \textit{Representação dos MACs do esquema limitado:} cada \acrshort{MAC} é armazenado como um vetor de coeficientes, com um byte por coeficiente, o que ocupa cerca de oito vezes o espaço estritamente necessário.
	\item \textit{Primitivas criptográficas:} a cifragem dos pacotes utiliza um fluxo de chave derivado da função SHAKE-256 e a função pseudoaleatória utiliza HMAC-SHA-256, ambas disponíveis na biblioteca padrão do Python; uma implementação de produção deveria empregar uma cifra autenticada padronizada.
\end{itemize}

\section{Validação}

O comportamento da implementação é verificado por uma suíte de testes automatizados que, além dos casos nominais, simula provedores desonestos: provedores que respondem incorretamente aos desafios, que substituem os dados após responder corretamente, que perdem pacotes ou etiquetas, que descartam os dados mantendo as etiquetas, que trocam etiquetas entre posições e que corrompem pares de pacotes com o mesmo padrão de erro. Os testes também verificam a classificação de polinômios irredutíveis e primitivos conhecidos, o uso único dos polinômios do esquema limitado, a realização de mais auditorias compactas do que o conjunto limitado permitiria sem que o polinômio secreto seja transmitido, a ausência de informação estrutural nos pacotes armazenados, a substituição de pacotes por símbolos de reparo com \acrshort{ESI}s nunca reutilizados, a recomposição dos \emph{chunks} após a saída de um provedor e a reconstrução do arquivo original após a Redistribuição.
